\documentclass[11pt]{article}
\usepackage[margin=1in]{geometry}
\usepackage{amsmath,amssymb,amsthm}
\usepackage{algorithm}
\usepackage{algpseudocode}
\usepackage{booktabs}
\usepackage{hyperref}

\newtheorem{definition}{Definition}
\newtheorem{theorem}{Theorem}
\newtheorem{lemma}[theorem]{Lemma}
\newtheorem{property}{Property}

\title{\textbf{Mathematical Specification, Modulated Gradient Dynamics, and Imbalance Theory of Focal Loss (ALGO-NN-17)}}
\author{Team Scaibu \\ \small Email: \texttt{jaydeep.9664920749@gmail.com} \\ \small Architecture and Core Engine Team}
\date{\today}

\begin{document}
\maketitle

\begin{abstract}
Dense object detection and extreme anomaly detection encounter severe class imbalance ratios (often exceeding $1000:1$). Standard cross-entropy gradients become completely overwhelmed by the cumulative weight of easily classified negative background examples. Focal Loss introduces a dynamically modulated loss functional $-\alpha_t (1 - p_t)^\gamma \log(p_t)$. This specification presents the formal mathematical foundations, derives the full gradient dynamics, proves polynomial background suppression, and provides complete pseudocode matching the reference implementation.
\end{abstract}

\section{Notation and Mathematical Preliminaries}
\label{sec:notation}

Let $z \in \mathbb{R}$ represent an unnormalized logit score, $p = \sigma(z) \in (0, 1)$ the predicted probability, and $y \in \{0, 1\}$ the binary target.

\subsection{Target-Aligned Probability and Modulating Weight}
\paragraph{Intuitive Meaning:} Focal loss scales the loss of each sample by $(1 - p_t)^\gamma$. When a prediction is highly confident and correct ($p_t \approx 0.99$), the factor $(1 - 0.99)^2 = 0.0001$ shrinks the loss by four orders of magnitude, preventing easy negatives from dominating optimization.

\paragraph{Formal Mathematical Definition:}
\begin{equation}
\label{eq:pt_def}
p_t = \begin{cases} p & \text{if } y = 1, \\ 1 - p & \text{if } y = 0, \end{cases} \quad \alpha_t = \begin{cases} \alpha & \text{if } y = 1, \\ 1 - \alpha & \text{if } y = 0. \end{cases}
\end{equation}
The modulated Focal Loss is defined as:
\begin{equation}
\label{eq:focal_def}
\mathcal{L}_{\mathrm{focal}}(p_t; \alpha, \gamma) = -\alpha_t (1 - p_t)^\gamma \log(p_t).
\end{equation}

\paragraph{Worked Numerical Example:}
Let $z = 2.19722 \implies p = 0.90, y = 1, \alpha = 0.25, \gamma = 2.0$:
$p_t = 0.90, \log(p_t) = -0.10536$.
Modulating factor: $(1 - 0.90)^2 = 0.01$.
Loss: $\mathcal{L} = -0.25 \times 0.01 \times (-0.10536) = 0.0002634$.

\subsection{Summary of Symbols}
\begin{itemize}
    \item $N \in \mathbb{N}$: Number of batch predictions.
    \item $\mathbf{z} \in \mathbb{R}^N$: Logit predictions.
    \item $\mathbf{y} \in \{0, 1\}^N$: Binary class labels.
    \item $\alpha \in (0, 1)$: Foreground balancing scalar.
    \item $\gamma \in \mathbb{R}_{\ge 0}$: Focusing parameter.
    \item $\mathbf{p} \in (0, 1)^N$: Predicted sigmoid probabilities.
    \item $\mathbf{g} \in \mathbb{R}^N$: Analytic gradients $dL/dz$.
\end{itemize}

\section{Problem Definition and Core Concepts}
\label{sec:problem_definition}

\subsection{Analytic Logit Gradient Derivation}
Differentiating Equation~\eqref{eq:focal_def} with respect to logit $z$:
\begin{equation}
\frac{\partial \mathcal{L}_{\mathrm{focal}}}{\partial z} = \alpha_t (1 - p_t)^\gamma \left[ \gamma p_t \log(p_t) + p_t - 1 \right] \cdot (2y - 1).
\end{equation}
When $\gamma = 0$ and $\alpha = 1$, this collapses to standard cross-entropy: $\frac{\partial \mathcal{L}}{\partial z} = p - y$.

\section{Algorithmic Specification}
\label{sec:algorithm}

Algorithm~\ref{alg:focal_loss} specifies the complete calculation routine matching \texttt{impl.py}.

\begin{algorithm}[H]
\caption{Focal Loss Forward and Gradient Evaluation with Stable Logits}
\label{alg:focal_loss}
\begin{algorithmic}[1]
\Require Logits $\mathbf{z} \in \mathbb{R}^N$, targets $\mathbf{y} \in \{0, 1\}^N$, $\alpha \in (0, 1)$, $\gamma \ge 0$, reduction $\mathrm{red} \in \{\text{mean}, \text{sum}, \text{none}\}$.
\Ensure Loss $\mathcal{L}$, probabilities $\mathbf{p} \in \mathbb{R}^N$, gradients $\mathbf{g} \in \mathbb{R}^N$.
\State Initialize $\mathbf{p} \gets [], \mathbf{l} \gets [], \mathbf{g} \gets []$
\For{$i = 1$ to $N$}
    \State $z \gets z_i, \quad y \gets y_i$
    \If{$z \ge 0$}
        \State $e_{-z} \gets \exp(-z)$
        \State $p_i \gets \frac{1.0}{1.0 + e_{-z}}, \quad \log p \gets -\log(1.0 + e_{-z}), \quad \log(1-p) \gets -z - \log(1.0 + e_{-z})$
    \Else
        \State $e_z \gets \exp(z)$
        \State $p_i \gets \frac{e_z}{1.0 + e_z}, \quad \log p \gets z - \log(1.0 + e_z), \quad \log(1-p) \gets -\log(1.0 + e_z)$
    \EndIf
    \If{$y = 1$}
        \State $p_t \gets p_i, \quad \log p_t \gets \log p, \quad \alpha_t \gets \alpha, \quad sign \gets 1.0$
    \Else
        \State $p_t \gets 1.0 - p_i, \quad \log p_t \gets \log(1-p), \quad \alpha_t \gets 1.0 - \alpha, \quad sign \gets -1.0$
    \EndIf
    \State $mod \gets (1.0 - p_t)^\gamma$
    \State $l_i \gets -\alpha_t \cdot mod \cdot \log p_t$
    \If{$\gamma = 0$}
        \State $g_i \gets \alpha_t \cdot (p_t - 1.0) \cdot sign$
    \Else
        \State $g_i \gets \alpha_t \cdot mod \cdot (\gamma \cdot p_t \cdot \log p_t + p_t - 1.0) \cdot sign$
    \EndIf
    \State Append $p_i$ to $\mathbf{p}$, $l_i$ to $\mathbf{l}$, $g_i$ to $\mathbf{g}$
\EndFor
\If{$\mathrm{red} = \text{mean}$}
    \State \Return $\frac{1}{N}\sum l_i, \quad \mathbf{p}, \quad \frac{1}{N}\mathbf{g}$
\ElsIf{$\mathrm{red} = \text{sum}$}
    \State \Return $\sum l_i, \quad \mathbf{p}, \quad \mathbf{g}$
\Else
    \State \Return $\mathbf{l}, \quad \mathbf{p}, \quad \mathbf{g}$
\EndIf
\end{algorithmic}
\end{algorithm}

\section{Theoretical Guarantees and Suppression Analysis}
\label{sec:guarantees}

\begin{theorem}[Polynomial Gradient Attenuation for Well-Classified Samples]
\label{thm:focal_polynomial_decay}
\textbf{Assumptions:} Let focusing exponent $\gamma \ge 1$, $\alpha_t \in (0, 1)$, and target probability $p_t \to 1.0$.
\textbf{Guarantee:} The absolute gradient magnitude $|\frac{\partial \mathcal{L}}{\partial z}|$ decays as $\mathcal{O}((1 - p_t)^{\gamma + 1})$, in contrast to standard cross-entropy which decays linearly as $\mathcal{O}(1 - p_t)$.
\textbf{Proof:}
As $p_t \to 1.0$, using the Taylor expansion $\log(p_t) = \log(1 - (1 - p_t)) = -(1 - p_t) + \mathcal{O}((1 - p_t)^2)$:
\begin{equation}
\gamma p_t \log(p_t) + p_t - 1 = -\gamma(1 - (1 - p_t))(1 - p_t) - (1 - p_t) = -(\gamma + 1)(1 - p_t) + \mathcal{O}((1 - p_t)^2).
\end{equation}
Multiplying by the modulating term $(1 - p_t)^\gamma$:
\begin{equation}
\left| \frac{\partial \mathcal{L}_{\mathrm{focal}}}{\partial z} \right| = \alpha_t (\gamma + 1) (1 - p_t)^{\gamma + 1} + \mathcal{O}((1 - p_t)^{\gamma + 2}).
\end{equation}
This establishes a polynomial suppression exponent of $\gamma + 1$.
\textbf{Limitations:} Output predictions $p$ deviate from true class posterior probabilities; score calibration is necessary for decision-theoretic risk thresholds.
\end{theorem}

\section{Complexity Analysis}
\label{sec:complexity}

\begin{itemize}
    \item \textbf{Time Complexity:} $\mathcal{O}(N)$ elementary operations.
    \item \textbf{Space Complexity:} $\mathcal{O}(N)$ gradient and loss memory buffers.
\end{itemize}

\section{Worked Numerical Example}
\label{sec:worked_example}

Let $z = 0.0 \implies p = 0.5$, $y = 1, \alpha = 0.25, \gamma = 2.0$:
\begin{enumerate}
    \item $p_t = 0.5, \log(p_t) = \log(0.5) = -0.693147$.
    \item Modulating factor: $(1 - 0.5)^2 = 0.25$.
    \item Loss: $\mathcal{L} = -0.25 \times 0.25 \times (-0.693147) = 0.043322$.
    \item Gradient: $mod \times (\gamma p_t \log p_t + p_t - 1) = 0.25 \times (2(0.5)(-0.693147) + 0.5 - 1) = 0.25 \times (-0.693147 - 0.5) = 0.25 \times (-1.193147) = -0.298287$.
    Scaled by $\alpha = 0.25 \implies g = 0.25 \times (-0.298287) = -0.074572$.
\end{enumerate}

\section{Numerical Considerations and Edge Cases}
\label{sec:numerical}

\begin{itemize}
    \item \textbf{High-Probability Clamping:} When $p_t > 1.0 - 10^{-15}$, numerical subtraction $(1 - p_t)$ is protected to prevent division by zero or underflow in power operations.
\end{itemize}

\begin{thebibliography}{9}
\bibitem{lin2017} T.-Y.~Lin et al., ``Focal Loss for Dense Object Detection,'' \emph{ICCV}, 2017.
\end{thebibliography}

\end{document}
